% Типографская обёртка условий вариантов 1–16.
% Смысл и данные совпадают с raw-submission.tex; исходный файл не редактируется.

% ВАРИАНТ 1
Цилиндр вписан в правильную четырёхугольную призму. Радиус основания и высота
цилиндра равны \(3\). Найдите площадь боковой поверхности призмы.

% ВАРИАНТ 2
Цилиндр вписан в правильную шестиугольную призму. Радиус основания цилиндра
равен \(\sqrt{3}\), а высота призмы равна \(2\). Найдите площадь боковой
поверхности призмы.

% ВАРИАНТ 3
Основанием пирамиды служит прямоугольник, одна боковая грань перпендикулярна
плоскости основания, а три другие боковые грани наклонены к плоскости основания
под углом \(60^\circ\). Высота пирамиды равна \(6\). Найдите объём пирамиды.

% ВАРИАНТ 4
Сторона основания правильной шестиугольной пирамиды равна \(3\), боковое ребро
равно \(6\). Найдите объём пирамиды.

% ВАРИАНТ 5
Высота конуса равна \(18\), а длина образующей равна \(30\). Найдите площадь
осевого сечения этого конуса.

% ВАРИАНТ 6
Диаметр основания конуса равен \(32\), а длина образующей равна \(20\). Найдите
площадь осевого сечения этого конуса.

% ВАРИАНТ 7
От треугольной призмы, объём которой равен \(120\), отсечена треугольная
пирамида плоскостью, проходящей через сторону одного основания и противоположную
вершину другого основания. Найдите объём оставшейся части.

% ВАРИАНТ 8
Объём треугольной пирамиды равен \(14\). Плоскость проходит через сторону
основания этой пирамиды и пересекает противоположное боковое ребро в точке,
делящей его в отношении \(2:5\), считая от вершины пирамиды. Найдите большую из
объёмов пирамид, на которые плоскость разбивает исходную пирамиду.

% ВАРИАНТ 9
Из единичного куба вырезана правильная четырёхугольная призма со стороной
основания \(0{,}4\) и боковым ребром \(1\). Найдите площадь поверхности
оставшейся части куба.

% ВАРИАНТ 10
Из единичного куба вырезана правильная четырёхугольная призма со стороной
основания \(0{,}6\) и боковым ребром \(1\). Найдите площадь поверхности
оставшейся части куба.

% ВАРИАНТ 11
В кубе \(ABCDA_1B_1C_1D_1\) найдите угол между прямыми \(DC_1\) и \(BD\).
Ответ дайте в градусах.

% ВАРИАНТ 12
В правильной треугольной призме \(ABCA_1B_1C_1\), все рёбра которой равны
\(2\), найдите угол между прямыми \(BB_1\) и \(AC_1\). Ответ дайте в градусах.

% ВАРИАНТ 13
Радиусы двух шаров равны \(7\) и \(24\). Найдите радиус шара, площадь
поверхности которого равна сумме площадей поверхностей двух данных шаров.

% ВАРИАНТ 14
Объём параллелепипеда \(ABCDA_1B_1C_1D_1\) равен \(60\). Найдите объём
треугольной пирамиды \(ACB_1D_1\).

% ВАРИАНТ 15
В прямоугольном параллелепипеде \(ABCDA_1B_1C_1D_1\) известно, что \(AB=9\),
\(BC=6\), \(AA_1=5\). Найдите объём многогранника, вершинами которого являются
точки \(A\), \(B\), \(C\), \(D\), \(A_1\), \(B_1\).

% ВАРИАНТ 16
В прямоугольном параллелепипеде \(ABCDA_1B_1C_1D_1\) известно, что \(AB=9\),
\(BC=8\), \(AA_1=6\). Найдите объём многогранника, вершинами которого являются
точки \(A\), \(B\), \(C\), \(B_1\).
